% [VERIFY] appendix -- src/configs/{training,model}.md, src/environment.md, experiments.md E06 setup, constraints.md.
% Implementation substrate routed here per methodology-discipline (body stays contribution-level).
\section{Implementation Details}
\label{app:impl}

\noindent\textbf{ImageNet training.}
Our ImageNet models follow the training recipe established for deep
batch-normalized networks. The image is resized with its shorter side sampled in
$[256, 480]$ for scale augmentation; a $224\times224$ crop is randomly sampled
from the image or its horizontal flip, with the per-pixel mean subtracted and
standard color augmentation applied. Batch normalization is adopted right after
each convolution and before activation. Weights are initialized with the MSRA
scheme, and all networks are trained from scratch. We use SGD with a mini-batch
size of $256$, momentum $0.9$, and weight decay $10^{-4}$. The learning rate
starts at $0.1$ and is divided by $10$ when the validation error plateaus, and
the models are trained for up to $60\times10^{4}$ iterations. We do not use
dropout. In testing, we report 10-crop results following standard
practice~\cite{alexnet2012} for the comparison tables; for the best single-model
results we adopt fully-convolutional multi-scale testing~\cite{vgg2015,overfeat2014,fcn2015}
at scales $\{224, 256, 384, 480, 640\}$ and average the scores. All activations
are ReLU~\cite{relu2010}.

\noindent\textbf{CIFAR-10 training.}
On CIFAR-10~\cite{cifar2009} we use a mini-batch size of $128$ on two GPUs, momentum $0.9$, and
weight decay $10^{-4}$, with MSRA initialization and batch normalization and no
dropout. Following standard practice~\cite{dsn2014}, we pad the $32\times32$ image by $4$ pixels
on each side and sample a random $32\times32$ crop from the padded image or its
horizontal flip; only the original $32\times32$ image is used at test time, with
per-pixel mean subtraction. The learning rate starts at $0.1$ and is divided by
$10$ at $32{,}000$ and $48{,}000$ iterations, with training terminated at
$64{,}000$ iterations on a $45$k/$5$k train/validation split. For the 110-layer
model only, we first warm up at a learning rate of $0.01$ for $\sim 400$
iterations until the training error is below $\sim 80\%$, then restore $0.1$.

\noindent\textbf{Detection fine-tuning.}
For detection we use the baseline Faster R-CNN system with 4-step alternating
training. The ResNet-101 backbone shares the full-image convolutional features
through the stage with total stride $16$; region-of-interest pooling is applied
before the final convolutional stage, which then acts as the per-region
equivalent of the fully-connected layers, with the classification and box
regression heads replaced for the target task. Batch-normalization statistics
are frozen during fine-tuning to reduce memory consumption, so batch
normalization behaves as an affine transform. The detection network is trained
on $8$ GPUs with a mini-batch of $8$ images for the region-proposal step and
$16$ for the detector step; the learning rate is $0.001$ for $240$k iterations
and then $0.0001$ for $80$k iterations.

\noindent\textbf{Environment.}
The original work does not commit to a specific framework in the manuscript; the
reference implementation of the era is a standard convolutional training
stack~\cite{caffe2014} with per-pixel mean subtraction and color/scale
augmentation data layers, with normalization following established
practice~\cite{efficientbackprop1998}. The ImageNet GPU count is not specified,
while the mini-batch size of $256$ implies a multi-GPU setup; CIFAR-10 uses two
GPUs and detection uses eight. Random seeds are not specified; the 110-layer
CIFAR model is reported as five runs with mean$\pm$std, so averaging over seeds
is performed for that depth.
